\documentclass[12pt]{article}
\usepackage[margin=1in]{geometry}
\usepackage{enumitem}
\usepackage{titlesec}
\usepackage{parskip}
\usepackage{graphicx}
\usepackage{xcolor}
\usepackage{hyperref}
\usepackage{fontenc}

\title{\textbf{North and Weingast 1989 RG Notes}\\
\large North, Douglass C., and Barry R. Weingast. ``Constitutions and Commitment: The Evolution of Institutions Governing Public Choice in Seventeenth-Century England'' [1989].\\
\textit{The Journal of Economic History} 49(4): 803-832.}
\author{}
\date{}

\begin{document}

\maketitle

Primary argument: development of wealth in England following 1688 Glorious Revolution occurred because the state adopted institutions allowing itself to credibly commit to property rights, most notably returning bonds; in turn, the wealthy were comfortable both acquiring more capital and lending it to the state, allowing for rapid devlopment

Can be largely connected to Huntington -- what political institutions exist as wealth develops? Definitely a different approach, though, as NW suggests that political development causes wealth

Self-enforcing constitution -- see Boix 1999 on parties changing electoral rules to favor themselves

``Our view also implies that the development of free markets must be accompanied by some credible restrictions on the state's ability to manipulate economic rules to the advantage of itself and its constituents. Successful economic performance, therefore, must be accompanied by institutions that limit economic intervention and allow private rights and markets to prevail in large segments of the economy.'' (North and Weingast, 1989, p. 808)

See Why the west extended the franchise for credibility A\&R 2000

\end{document}
